%%%%%%%%%%%%%%%%%%%%%%%%%%%%%%%%%%%%%%%%%%%%%%%%%%%%%%%%%%%%%%%%%%%%%%
% Lớp 8 | Chương III. Tứ giác | Hình thang cân
% Dạng 8H3?11-2: Dấu hiệu nhận biết hình thang cân
%%%%%%%%%%%%%%%%%%%%%%%%%%%%%%%%%%%%%%%%%%%%%%%%%%%%%%%%%%%%%%%%%%%%%%


% nguồn: De_AI_Toan8_cuoi_hoc_ki_i_1 (nhập 21/11/2026)
\begin{bt}%[8H3H11-2]%[ma:d28c9ba6]
	Cho tam giác $ABC$ cân tại $A$. Trên cạnh $AB$ lấy điểm $D$, trên cạnh $AC$ lấy điểm $E$ sao cho $AD = AE$.
a) Chứng minh rằng $DE \parallel BC$.
b) Chứng minh rằng tứ giác $BDEC$ là hình thang cân.
	\loigiai{
		a) Vì tam giác $ABC$ cân tại $A$ nên $\widehat{B} = \dfrac{180^\circ - \widehat{A}}{2}$. \\
Vì $AD = AE$ nên tam giác $ADE$ cân tại $A$, suy ra $\widehat{ADE} = \dfrac{180^\circ - \widehat{A}}{2}$. \\
Do đó $\widehat{ADE} = \widehat{B}$. Mà hai góc này ở vị trí đồng vị nên $DE \parallel BC$. \\
b) Xét tứ giác $BDEC$ có $DE \parallel BC$ nên $BDEC$ là hình thang. \\
Lại có $\widehat{B} = \widehat{C}$ (vì tam giác $ABC$ cân tại $A$), nên hình thang $BDEC$ có hai góc kề một đáy bằng nhau, do đó $BDEC$ là hình thang cân.
	}
\end{bt}

% nguồn: De_AI_Toan8_cuoi_hoc_ki_i_1 (nhập 21/11/2026)
\begin{bt}%[8H3H11-2]%[ma:68629572]
	Cho tam giác $ABC$ cân tại $A$. Các đường phân giác của góc $B$ và góc $C$ lần lượt cắt cạnh $AC$ tại $E$ và cắt cạnh $AB$ tại $D$.
a) Chứng minh tam giác $ABD$ bằng tam giác $ACE$, từ đó suy ra $AD = AE$.
b) Chứng minh tứ giác $BDEC$ là hình thang cân.
	\loigiai{
		a) Vì tam giác $ABC$ cân tại $A$ nên $AB = AC$ và $\widehat{ABC} = \widehat{ACB}$. \\
Vì $BE$, $CD$ là phân giác của $\widehat{ABC}$, $\widehat{ACB}$ nên $\widehat{ABE} = \widehat{ACD} = \dfrac{1}{2}\widehat{ABC}$. \\
Xét $\triangle ABE$ và $\triangle ACD$ có: $\widehat{A}$ chung, $AB = AC$, $\widehat{ABE} = \widehat{ACD}$. Suy ra $\triangle ABE = \triangle ACD$ (g.c.g), do đó $AE = AD$. \\
b) Vì $AD = AE$ nên tam giác $ADE$ cân tại $A$, suy ra $\widehat{ADE} = \dfrac{180^\circ - \widehat{A}}{2} = \widehat{ABC}$, do đó $DE \parallel BC$. \\
Hình thang $BDEC$ ($DE \parallel BC$) có $\widehat{B} = \widehat{C}$ nên là hình thang cân.
	}
\end{bt}

% nguồn: De_AI_Toan8_cuoi_hoc_ki_i_1 (nhập 21/11/2026)
\begin{bt}%[8H3H11-2]%[ma:9312d97f]
	Cho hình thang $ABCD$ ($AB \parallel CD$) có hai đường chéo $AC = BD$ và $\widehat{ADC} = 72^\circ$.
a) Tứ giác $ABCD$ là hình gì? Vì sao?
b) Tính số đo các góc còn lại của tứ giác $ABCD$.
	\loigiai{
		a) Hình thang $ABCD$ ($AB \parallel CD$) có hai đường chéo bằng nhau ($AC = BD$) nên theo dấu hiệu nhận biết, $ABCD$ là hình thang cân. \\
b) Vì $ABCD$ là hình thang cân ($AB \parallel CD$) nên: \\
$\widehat{BCD} = \widehat{ADC} = 72^\circ$. \\
$\widehat{DAB} = 180^\circ - \widehat{ADC} = 180^\circ - 72^\circ = 108^\circ$. \\
$\widehat{ABC} = \widehat{DAB} = 108^\circ$.
	}
\end{bt}
